\section{DecMed}

DecMed merupakan sebuah kerangka kerja sistem manajemen Rekam Medis Elektronik (RME) berbasis teknologi Distributed Ledger Technology (DLT) IOTA yang dirancang untuk meningkatkan keamanan data, khususnya aspek kerahasiaan dan integritas. Sistem ini mengintegrasikan IOTA sebagai infrastruktur ledger, InterPlanetary File System (IPFS) sebagai media penyimpanan data, serta Proxy Re-Encryption sebagai mekanisme enkripsi dan delegasi akses. DecMed menerapkan model kontrol akses yang berpusat pada pasien, di mana setiap akses terhadap data RME harus memperoleh persetujuan aktif dari pasien, disertai kemampuan untuk mengatur durasi akses dan membatasi bagian data yang dapat diakses. Selain itu, DecMed mendukung pengelolaan akun pasien dan tenaga kesehatan yang terintegrasi dengan mekanisme kontrol akses, serta telah divalidasi mampu memenuhi kebutuhan manajemen RME di Indonesia.

\subsection{Arsitektur DecMed}

Arsitektur DecMed terdiri atas tiga aplikasi klien, yaitu klien pasien, klien fasilitas pelayanan kesehatan (digunakan oleh tenaga kesehatan, petugas rekam medis, dan administrator), serta klien Kementerian Kesehatan. Sistem ini juga mencakup komponen pendukung berupa Proxy Re-Encryption Server untuk proses delegasi akses data, Gas Station Server untuk mensponsori transaksi pada jaringan IOTA, Redis untuk menyimpan data pendukung mekanisme proxy re-encryption, dan IPFS Cluster sebagai media penyimpanan data rekam medis elektronik yang telah dienkripsi. Dalam mekanisme penyimpanannya, DecMed menggunakan IOTA untuk menyimpan data administratif pasien, data administratif tenaga kesehatan, serta metadata rekam medis berupa kunci enkripsi dan CID yang mengarah ke data pada IPFS. Sementara itu, bagian klinis rekam medis disimpan pada IPFS dalam keadaan terenkripsi, sehingga data rekam medis tersebar antara IOTA dan IPFS sesuai fungsinya. Gambar \ref{fig:arsitektur-decmed} menunjukkan arsitektur sistem DecMed secara keseluruhan.

\begin{figure}[ht]
    \centering
    \includegraphics[width=\textwidth]{chapter-2/decmed_archi.png}
    \caption{Gambar Arsitektur Decmed}
    \label{fig:arsitektur-decmed}
\end{figure}

\subsubsection{InterPlanetary File System (IPFS)}
InterPlanetary File System (IPFS) merupakan sistem berkas terdistribusi berbasis peer-to-peer (P2P) yang dirancang untuk menyediakan mekanisme penyimpanan dan distribusi data secara terdesentralisasi. Berbeda dengan model penyimpanan konvensional yang bergantung pada server pusat, IPFS memungkinkan setiap node dalam jaringan berperan sebagai penyedia maupun konsumen data. Dengan pendekatan tersebut, IPFS tidak memiliki single point of failure dan tidak mengharuskan setiap node saling mempercayai satu sama lain untuk melakukan pertukaran data. Selain itu, IPFS menggunakan mekanisme content addressing, yaitu setiap objek diidentifikasi berdasarkan nilai hash dari isi datanya, bukan berdasarkan lokasi penyimpanannya. Pendekatan ini memungkinkan data yang identik memiliki identitas yang sama sehingga dapat dilakukan deduplikasi secara otomatis.

Secara konseptual, IPFS menggabungkan berbagai konsep dari sistem terdistribusi yang telah terbukti berhasil, seperti Distributed Hash Table (DHT), BitTorrent, Git, serta Self-Certified File System (SFS). Integrasi berbagai teknologi tersebut menghasilkan sebuah sistem yang mampu mendukung penyimpanan objek secara terdistribusi, distribusi data yang efisien, mekanisme pengalamatan berbasis konten, serta pengelolaan versi data menggunakan struktur Merkle Directed Acyclic Graph (Merkle DAG). Struktur ini memungkinkan IPFS tidak hanya berfungsi sebagai sistem penyimpanan berkas, tetapi juga sebagai fondasi bagi berbagai aplikasi lain seperti blockchain, sistem version control, maupun penyimpanan data permanen.

\paragraph{Keamanan IPFS}
Keamanan pada IPFS dibangun melalui beberapa mekanisme yang saling melengkapi. Mekanisme pertama adalah content addressing, yaitu setiap objek diidentifikasi menggunakan hash kriptografi dari isi datanya. Ketika sebuah objek diterima dari node lain, IPFS secara otomatis menghitung kembali nilai hash objek tersebut dan membandingkannya dengan hash yang diminta. Apabila kedua nilai tersebut tidak sama, maka objek dianggap telah mengalami perubahan atau kerusakan sehingga dapat ditolak. Pendekatan ini memberikan jaminan integritas data tanpa memerlukan kepercayaan terhadap node pengirim.

Mekanisme kedua adalah penggunaan Merkle DAG. Karena setiap objek menyimpan hash dari objek lain yang dirujuk, perubahan pada satu objek akan mengubah hash seluruh jalur yang mengarah ke objek tersebut. Dengan demikian, setiap modifikasi terhadap data dapat dideteksi secara otomatis. Selain mendukung verifikasi integritas, struktur Merkle DAG juga memungkinkan deduplikasi data karena objek dengan isi yang identik hanya perlu disimpan satu kali.

Pada tingkat identitas, setiap node memiliki pasangan public key dan private key yang digunakan untuk membentuk Node ID. Ketika dua node membangun koneksi, keduanya saling memverifikasi kesesuaian antara public key dan Node ID sehingga identitas node dapat dipastikan sebelum komunikasi dilakukan. Selain itu, desain IPFS mengadopsi pendekatan dari S/Kademlia untuk meningkatkan ketahanan terhadap serangan seperti Sybil Attack melalui mekanisme pembangkitan identitas berbasis kriptografi.

IPFS juga menyediakan dukungan terhadap kriptografi pada tingkat objek. Suatu objek dapat dienkripsi ataupun ditandatangani secara digital sebelum dipublikasikan ke jaringan. Ketika objek dienkripsi, hash objek akan berubah sehingga menghasilkan identitas baru yang berbeda dari objek asli. Selain itu, tautan menuju objek terenkripsi juga ikut terlindungi sehingga objek tidak dapat diakses maupun ditelusuri tanpa memiliki kunci dekripsi yang sesuai. IPFS secara otomatis melakukan verifikasi terhadap tanda tangan digital serta mendukung proses dekripsi menggunakan kumpulan kunci yang dimiliki pengguna.

Meskipun demikian, IPFS tidak secara otomatis menjamin kerahasiaan data yang disimpan. Objek yang dipublikasikan ke jaringan pada dasarnya dapat diakses oleh siapa saja yang mengetahui nilai hash atau Content Identifier (CID) dari objek tersebut. Oleh karena itu, data yang bersifat sensitif perlu dienkripsi terlebih dahulu sebelum disimpan di IPFS sehingga hanya pihak yang memiliki kunci dekripsi yang dapat mengakses isi data tersebut.

\subsection{Audit Trail DecMed}

Salah satu tujuan arsitektur DecMed adalah meningkatkan auditability melalui pemanfaatan Distributed Ledger Technology (DLT) sebagai media penyimpanan yang bersifat tamper-evident. Audit trail merupakan salah satu kebutuhan untuk mendukung akuntabilitas, kepatuhan terhadap regulasi, serta kesiapan forensik, terutama pada lingkungan layanan kesehatan yang melibatkan banyak institusi.

Implementasi auditability pada DecMed dicapai melalui pencatatan aktivitas kontrol akses ke dalam jaringan IOTA. Setiap perubahan hak akses, seperti pemberian (\textit{grant}) dan pencabutan (\textit{revoke}) akses terhadap data RME, direpresentasikan sebagai transaksi yang tersimpan secara permanen pada ledger. Selain itu, metadata rekam medis dan capability yang digunakan dalam mekanisme Capability-Based Access Control juga disimpan pada IOTA sehingga keberadaan dan validitas hak akses dapat diverifikasi melalui smart contract. Dengan memanfaatkan sifat immutable dari DLT, catatan tersebut tidak dapat diubah secara sepihak sehingga meningkatkan integritas bukti akses.

% Meskipun demikian, implementasi audit trail pada DecMed masih terbatas pada pencatatan transaksi yang berkaitan dengan mekanisme kontrol akses dan metadata sistem. Paper tidak menjelaskan adanya komponen audit trail yang berdiri sendiri, seperti mekanisme pengumpulan audit event dari seluruh komponen sistem, struktur audit record yang terstandarisasi, penyimpanan log khusus, mekanisme retensi, indeks dan pencarian log, maupun fasilitas analisis audit. Dengan demikian, auditability pada DecMed terutama diperoleh dari karakteristik immutable jaringan IOTA, sementara infrastruktur audit trail yang komprehensif sebagaimana direkomendasikan dalam standar seperti NIST SP 800-92 atau NIST SP 800-53 belum menjadi fokus pembahasan dalam penelitian tersebut.

\subsection{Komponen dan Struktur Rekam Medis Elektronik}

Sistem informasi rumah sakit sering kali menghadapi masalah kualitas dalam pencatatan jejak audit (\textit{audit trails}). Analisis terhadap sistem informasi rumah sakit menunjukkan bahwa mekanisme audit yang ada sering kali tidak memadai untuk merekonstruksi insiden keamanan secara utuh atau menjamin integritas data jangka panjang \cite{6}.

Selain itu, integrasi teknologi baru seperti \textit{blockchain} dalam layanan kesehatan telah diidentifikasi memiliki potensi keuntungan sekaligus ancaman baru. Meskipun menjanjikan desentralisasi, penerapannya harus mempertimbangkan interoperabilitas dan skalabilitas dalam lingkungan medis yang kompleks \cite{1}.



\subsection{Audit Trail DecMed}

DecMed sendiri sudah membuat sebuah audit trail sederhana dalam bentuk log berupa patient\_access\_log. Patient\_access\_log berisikan riwayat pemberian akses CRUD rekam medis pasien kepada pegawai fasyankes. Log tersebut berisikan Timestamp, Nama Fasyankes, NAMA Petugas, jenis akses (read/update), tipe data akses (medical/administrative). Struct dari PatienAccessLog berbentuk demikian.

% \begin{lstlisting}[language=Rust, caption={Struktur Data Patient Access Log}]
% public struct PatientAccessLog has copy, drop, store {
%     access_data_type: vector<HospitalPersonnelAccessDataType>,
%     access_type: HospitalPersonnelAccessType,
%     date: String,
%     exp_dur: u64,
%     hospital_metadata: HospitalMetadata,
%     hospital_personnel_address: address,
%     hospital_personnel_metadata: String,
%     index: u64,
%     is_revoked: bool,
% }
% \end{lstlisting}

% Proses generasi log dilakukan melalui 